% =====================================================================
%  User Manual: how students and admins use the web application.
%  Build: ./scripts/build_docs.sh   (or compile in docs/ with pdflatex)
% =====================================================================
\documentclass[11pt,a4paper]{article}
\usepackage{campusdoc}

\begin{document}

\makedoctitle{User Manual}{How students and administrators use Campus Events}{CEMP-UM-01}

\tableofcontents
\newpage

% =====================================================================
\section{About this manual}

Campus Events is a website where \textbf{students} find campus events and sign up for them, and
\textbf{administrators} (admins) manage events, venues, participant lists and reports. This manual
explains every screen from the point of view of the person using it. It does not assume any
programming knowledge.

\begin{itemize}
  \item Section~\ref{sec:start} covers getting started: opening the site, creating an account and logging in.
  \item Section~\ref{sec:student} is for \textbf{students}.
  \item Section~\ref{sec:admin} is for \textbf{admins}.
  \item Section~\ref{sec:messages} lists every message the site can show and what to do about it.
\end{itemize}

\begin{note}[Who installs the site?]
This manual assumes the site is already running. Installing it is described in the
\emph{Installation \& Operations Guide}. On a computer where it has been set up, the site is
usually started with \code{./scripts/run.sh} and opened at \url{http://127.0.0.1:5000}.
\end{note}

\subsection{Roles at a glance}

\begin{table}[H]
\centering
\small
\begin{tabularx}{\linewidth}{@{}X c c c@{}}
\toprule
\textbf{What you can do} & \textbf{Visitor} & \textbf{Student} & \textbf{Admin} \\
\midrule
See the home page and browse / search upcoming events & \checkmark & \checkmark & \checkmark \\
Open an event's detail page                            & \checkmark & \checkmark & \checkmark \\
Register for an event and cancel a registration        &            & \checkmark &            \\
See \menu{My registrations}                             &            & \checkmark &            \\
Open the admin dashboard                                &            &            & \checkmark \\
Create, edit and delete venues and events               &            &            & \checkmark \\
See participant lists and export them as CSV            &            &            & \checkmark \\
See reports                                             &            &            & \checkmark \\
\bottomrule
\end{tabularx}
\caption{What each type of user can do.}
\end{table}

A \textbf{visitor} is anyone who has not logged in. Anyone can create a \textbf{student} account.
\textbf{Admin} accounts cannot be created through the website; they are made by whoever installs the
site (Installation Guide, Section ``Admin account'').

% =====================================================================
\section{Getting started}
\label{sec:start}

\subsection{The navigation bar}

Every page has the same blue bar at the top. Its links change depending on who is logged in:

\begin{table}[H]
\centering
\small
\begin{tabularx}{\linewidth}{@{}l X@{}}
\toprule
\textbf{You are\dots} & \textbf{Links in the top bar} \\
\midrule
Not logged in  & \menu{Home}, \menu{Events}, \menu{Login}, \menu{Register} \\
A student      & \menu{Home}, \menu{Events}, \menu{My registrations}, ``Hi, \emph{your name}'', \menu{Logout} \\
An admin       & \menu{Home}, \menu{Events}, \menu{Admin}, ``Hi, \emph{your name}'', \menu{Logout} \\
\bottomrule
\end{tabularx}
\caption{Navigation bar links.}
\end{table}

Short messages (for example ``You are registered for Music Night.'') appear in a box just under
the bar after you do something. A \textcolor{okgreen}{green} edge means it worked; a
\textcolor{errred}{red} edge means something was refused, with the reason.

\subsection{The home page}

The home page (\menu{Home}) welcomes you and shows how many events, venues and students are in the
system.

\subsection{Creating a student account}

\begin{enumerate}
  \item Click \menu{Register} in the top bar.
  \item Fill in \textbf{Full name}, \textbf{Email} and \textbf{Password} (at least 6 characters).
  \item Click \button{Register}. If everything is valid you are taken to the login page with the
        message ``Account created. Please log in.''
\end{enumerate}

The site refuses the form if the name is empty, the email has no \code{@}, the password is shorter
than 6 characters, or an account with that email already exists. Emails are not case-sensitive:
\code{Asha@Campus.edu} and \code{asha@campus.edu} are the same account.

\subsection{Logging in and out}

Click \menu{Login}, enter your email and password and click \button{Log in}. You are taken to the
events list and greeted by name. If the email or password is wrong you see
``Invalid email or password.'' (The site deliberately does not say which one was wrong, so nobody
can use the login form to find out which emails have accounts.)

To log out, click \menu{Logout} in the top bar.

\begin{tip}
If you open a page that needs a login (such as \menu{My registrations}) while logged out, the
site sends you to the login page with ``Please log in to see that page.''
\end{tip}

% =====================================================================
\section{For students}
\label{sec:student}

\subsection{Browsing upcoming events}

Click \menu{Events}. The page \textbf{Upcoming events} shows every event that has not started yet,
earliest first, as a grid of cards. Each card shows the category, title, date and time, venue and
how many \textbf{seats are left}. Past events are never listed.

\begin{figure}[H]
\centering
\shot{search_page.png}{0.92\linewidth}
\caption{The Upcoming events page with the search form at the top.}
\label{fig:um-search}
\end{figure}

\subsection{Searching and filtering}

The search form at the top of the events page has five fields. Fill in any combination; empty
fields are ignored.

\begin{table}[H]
\centering
\small
\begin{tabularx}{\linewidth}{@{}l X@{}}
\toprule
\textbf{Field} & \textbf{What it does} \\
\midrule
Keyword  & Finds events whose \textbf{title or description} contains the text, e.g.\ \code{python}.
           Upper and lower case do not matter. \\
Category & One of \emph{Technical, Cultural, Sports, Workshop, Seminar}. \\
Venue    & Only events at that venue. \\
From     & Only events starting on or after this date. \\
To       & Only events starting on or before this date (the whole day is included). \\
\bottomrule
\end{tabularx}
\caption{Search fields.}
\end{table}

Click \button{Search} to apply the filters. The number of matching events is shown above the
cards. Click \textbf{Clear} to remove all filters.

\begin{tip}[Bookmark a search]
The filters appear in the page address, for example
\route{/events/?q=python&category=Workshop}. You can bookmark this address or send it to a friend
to repeat the same search later.
\end{tip}

\subsection{Viewing an event}

Click an event's title to open its detail page. It shows the category, title, \textbf{When}
(start and end), \textbf{Where} (venue and location), the \textbf{Organiser}, \textbf{Seats}
(how many are left out of the total) and the full description.

\begin{figure}[H]
\centering
\shot{detail_page.png}{0.92\linewidth}
\caption{An event detail page with the \textsf{Register} button.}
\end{figure}

The bottom of the page shows one of the following, depending on the situation:

\begin{table}[H]
\centering
\small
\begin{tabularx}{\linewidth}{@{}l X@{}}
\toprule
\textbf{Situation} & \textbf{What you see} \\
\midrule
You are not logged in                    & ``Log in to register for this event.'' \\
You are an admin                         & A note that admins do not register for events \\
You are registered, event not started    & \button{Cancel registration} \\
You are registered, event has started    & A note that you are registered \\
Event has started                        & A note that registration is closed \\
No seats left                            & A note that the event is full \\
Otherwise                                & \button{Register} \\
\bottomrule
\end{tabularx}
\caption{The action area of the event detail page.}
\end{table}

\subsection{Registering for an event}

\begin{enumerate}
  \item Log in as a student.
  \item Open the event and click \button{Register}.
  \item You stay on the event page and see ``You are registered for \emph{event title}.''
        The seat count goes down by one and the button changes to \textsf{Cancel registration}.
\end{enumerate}

The site checks the following, in this order, and tells you which rule stopped the registration:

\begin{enumerate}
  \item admins cannot register for events;
  \item the event must not have started yet;
  \item you must not already be registered for it;
  \item there must be at least one seat left.
\end{enumerate}

\subsection{My registrations}

Click \menu{My registrations} in the top bar to see every upcoming event you are registered for,
with its category, time, venue and when you registered. Each row has a red \textsf{Cancel} button.

\begin{figure}[H]
\centering
\shot{my_registrations_page.png}{0.92\linewidth}
\caption{The My registrations page.}
\end{figure}

\subsection{Cancelling a registration}

Click \textsf{Cancel registration} on the event page, or \textsf{Cancel} on the
\menu{My registrations} page. You return to the page you came from with
``Your registration for \emph{event title} was cancelled.'' Your seat becomes free for someone else.
A registration cannot be cancelled once the event has started.

% =====================================================================
\section{For administrators}
\label{sec:admin}

Log in with an admin account and click \menu{Admin} in the top bar. Every admin page is protected:
a student who types an admin address sees ``Error 403: You are not allowed to view this page.'', and
a visitor is sent to the login page.

\subsection{The dashboard}

The \textbf{Admin dashboard} shows four counters: \emph{Events} (with how many are upcoming),
\emph{Venues}, \emph{Students} and \emph{Registrations}. Below them are shortcut buttons:
\button{Manage events}, \button{Manage venues}, \button{Reports}, \textsf{+ New event} and
\textsf{+ New venue}.

\subsection{Managing venues}

\menu{Manage venues} lists every venue with its capacity, location and number of events, plus
\textsf{Edit} and \textsf{Delete} actions.

\paragraph{Create a venue.} Click \textsf{+ New venue}, fill in \textbf{Name}, \textbf{Location} and
\textbf{Capacity} (a whole number above 0), and click \button{Save}.

\paragraph{Edit a venue.} Click \textsf{Edit}, change the fields and click \button{Save}. The
capacity cannot be made smaller than the largest event already booked there.

\paragraph{Delete a venue.} Click \textsf{Delete} and confirm the browser's ``Delete this venue?''
question. A venue that still has events cannot be deleted; move or delete those events first.

\subsection{Managing events}

\menu{Manage events} lists \emph{all} events, past and upcoming, with category, start time, venue,
number registered and the actions \textsf{Participants}, \textsf{Edit} and \textsf{Delete}.

\paragraph{Create an event.} Click \textsf{+ New event} and fill in the form:

\begin{table}[H]
\centering
\small
\begin{tabularx}{\linewidth}{@{}l X@{}}
\toprule
\textbf{Field} & \textbf{Rules} \\
\midrule
Title        & Required. \\
Description  & Optional. Shown on the event page and searched by the keyword filter. \\
Category     & Required; one of the five categories. \\
Venue        & Required. The dropdown shows each venue's capacity, e.g.\ ``Sports Complex (holds 300)''. \\
Starts, Ends & Required date and time. The event must end after it starts. \\
Capacity     & Whole number above 0, \textbf{not more than the venue's capacity}. When editing, not less
               than the number of students already registered. \\
\bottomrule
\end{tabularx}
\caption{Event form fields.}
\end{table}

Click \button{Save}. If something is wrong, the form is shown again with your values and an error
message at the top, so nothing you typed is lost.

\paragraph{Venue clashes.} A venue cannot hold two events at the same time. If the new times overlap
another event at the same venue, the form is refused with a message such as:

\begin{quote}\small
Sports Complex is already booked for `Poetry Slam' (2026-11-10 13:00 to 2026-11-10 16:00).
\end{quote}

Events that only touch, for example one ending at 11:00 and the next starting at 11:00, are
allowed. When you edit an event, it is never reported as clashing with itself.

\begin{figure}[H]
\centering
\shot{clash_error_page.png}{0.92\linewidth}
\caption{A new event that overlaps another booking at the same venue is rejected.}
\end{figure}

\paragraph{Delete an event.} Click \textsf{Delete} and confirm ``Delete this event and all its
registrations?''. All registrations for that event are removed automatically.

\begin{warning}
Deleting an event or venue cannot be undone.
\end{warning}

\subsection{Participants and CSV export}

Click \textsf{Participants} next to an event to see the numbered list of registered students: name,
email and registration time, earliest first. Click \button{Export CSV} to download the same list as
\code{event_<id>_participants.csv}, which opens in Excel, LibreOffice Calc or Google Sheets. The file
has the columns \emph{Name}, \emph{Email} and \emph{Registered at}.

\subsection{Reports}

Click \button{Reports} on the dashboard. The page has three tables.

\begin{description}
  \item[Popular categories] For each category: number of events and total registrations, most
        popular first.
  \item[Venue utilisation] For each venue: events held, registrations, \emph{seats offered} (the
        total capacity of its events) and \emph{utilisation}, which is registrations as a
        percentage of seats offered. Venues with no events are listed with 0\%.
  \item[Registrations per event] Every event with its category, start, number registered,
        capacity and \emph{fill rate} (registered as a percentage of capacity), shown as a number
        and as a small bar. Most registered first.
\end{description}

\begin{figure}[H]
\centering
\shot{reports_overview.png}{0.92\linewidth}
\caption{Reports: popular categories and venue utilisation.}
\end{figure}

\begin{figure}[H]
\centering
\shot{reports_per_event.png}{0.92\linewidth}
\caption{Reports: registrations per event with fill-rate bars.}
\end{figure}

% =====================================================================
\section{Messages and what they mean}
\label{sec:messages}

\begin{xltabular}{\linewidth}{@{}>{\raggedright\arraybackslash}p{0.42\linewidth} X@{}}
\toprule
\textbf{Message} & \textbf{Meaning and what to do} \\
\midrule
\endhead
\multicolumn{2}{@{}l}{\textit{Accounts}} \\
Name is required. & Type your name. \\
Please enter a valid email address. & The email must contain \code{@}. \\
Password must be at least 6 characters. & Choose a longer password. \\
An account with that email already exists. & Log in instead, or use another email. \\
Invalid email or password. & Check both and try again. \\
Please log in to see that page. & Log in; you will be able to open the page afterwards. \\
\midrule
\multicolumn{2}{@{}l}{\textit{Registrations (students)}} \\
Admins cannot register for events. & Use a student account. \\
This event has already started. & Registration closes when the event starts. \\
You are already registered for this event. & Nothing to do; see \menu{My registrations}. \\
Sorry, this event is full. & All seats are taken. Check again later in case someone cancels. \\
You cannot cancel after the event has started. & Cancellation closes when the event starts. \\
You were not registered for this event. & There was no registration to cancel. \\
\midrule
\multicolumn{2}{@{}l}{\textit{Events and venues (admins)}} \\
Title is required. / Please choose a category. / Please choose a venue. & Fill in the missing field. \\
Please enter a valid start and end time. & Pick both a date and a time. \\
The event must end after it starts. & Fix the start or end time. \\
Capacity must be a whole number greater than 0. & Enter e.g.\ \code{50}. \\
Capacity (N) cannot be more than the capacity of \emph{venue} (M). & Lower the capacity or choose a bigger venue. \\
N students are already registered, so capacity cannot be lower. & Keep the capacity at least N. \\
\emph{Venue} is already booked for `\emph{event}' (\emph{times}). & Choose another time or venue. \\
Venue name is required. / Location is required. & Fill in the missing field. \\
An event here has capacity N, so the venue cannot be smaller. & Keep the venue capacity at least N. \\
Cannot delete: N event(s) use this venue. & Move or delete those events first. \\
\midrule
\multicolumn{2}{@{}l}{\textit{Error pages}} \\
Error 403: You are not allowed to view this page. & The page is for admins only. \\
Error 404: That page does not exist. & The address is wrong, or the event was deleted. \\
\bottomrule
\end{xltabular}

% =====================================================================
\section{Frequently asked questions}

\paragraph{Can I register a friend?} No. Each student registers themselves from their own account.

\paragraph{Why can't I see an event I registered for in My registrations?} The page only shows
events that have not started yet.

\paragraph{Why does an event say ``0 seats left'' but I can still see it?} Full events stay visible
so you know they exist; only the \textsf{Register} button is replaced.

\paragraph{How do I become an admin?} Admin accounts are created by whoever runs the site, using
\code{./scripts/setup.sh} or the \code{create-admin} command (see the Installation Guide).

\paragraph{Does the site work on a phone?} Yes. The layout adapts to narrow screens, and wide tables
can be scrolled sideways.

\paragraph{Demo accounts.} After the demo data is loaded with \code{./scripts/seed.sh}, you can log
in as the admin \code{admin@campus.local} / \code{admin123} or as any student from
\code{student1@campus.local} to \code{student100@campus.local} with password \code{student123}.
These are for demonstrations only.

\end{document}
